\section{Background and Related Work}\label{sec:background}
% Principal structure:
%
% 2.1 Provenance models — PROV-O \cite{prov-o2013} and its lineage (OPM
%     \cite{opm2010}); PAV for authoring/versioning \cite{pav2013}.
%     Position: aiprov extends, never forks (subclassing prov: terms).
% 2.2 Contributor attribution — CRediT \cite{credit2019,credit-niso2022}:
%     role taxonomies name WHO contributed WHAT KIND of work, but at
%     manuscript granularity and human-only. aiprov records per-activity,
%     per-claim, and admits AI + tool agents.
% 2.3 AI transparency artefacts — model cards \cite{modelcards2019},
%     datasheets \cite{datasheets2021}: document the MODEL or the DATA,
%     not the collaboration that produced a specific artefact.
% 2.4 Research packaging — RO-Crate \cite{rocrate2022}: complementary;
%     aiprov graphs can travel inside a crate.
% 2.5 The trust gap — hallucinated citations \cite{ai-hallucination2023}
%     preview the systemic problem treated in §\ref{sec:dilution}: the
%     ladder makes existence machine-checkable and support a judgement
%     with an explicit grantor.
% 2.6 Lineage and neighbors — direct ancestors: Obscurity-Is-Dead
%     \cite{obscurity-is-dead} (transcript-as-artifact; the labels
%     repo-vendored / lit-read / unverified-external that this ladder
%     canonicalizes) and F(AI)2R \cite{fair2r-repo} (two-pass model,
%     no-parentless-claim). Adjacent DLR work extends PROV-O with
%     uncertainty quantification for engineering traceability
%     \cite{uncertainty-prov2026}: complementary axes — aiprov records
%     who/how work happened, uncertainty-aware provenance records how
%     confident results are; a combined profile is future work (§8).
%     A third root is industrial semantic modeling: the author
%     co-authored the Asset Administration Shell Part 1 specification
%     \cite{aas-part1-v30rc02} (now IDTA-01001 \cite{aas-part1-idta})
%     — where partners exchange information across a value chain,
%     semantics must be formal and machine-readable or interoperability
%     fails. aiprov applies that same discipline to a new exchange
%     relationship: the one between humans, AI agents, and future
%     auditors of their joint work.
% 2.7 Metadata-community context — the Helmholtz Metadata Collaboration
%     (HMC) \cite{hmc-conference2025} pushes FAIR metadata practice
%     across research infrastructures; a provenance graph is metadata
%     ABOUT WORK, so aiprov can be read as extending that programme
%     from data description to collaboration description.
